\section{Implementation Details}
\label{sec:impl}

Our system is trained in two stages on NVIDIA H200 GPUs. Stage~1 pretrains
the trajectory encoder on StreamGaze\_v2 with trajectory-prediction and
patch-importance objectives (\S\ref{sec:impl-stage1}). Stage~2 jointly
fine-tunes LoRA adapters on the VLM together with the trajectory encoder
under a token-merging bottleneck (\S\ref{sec:impl-stage2}). For
self-containedness we first summarize the trajectory-encoder architecture
and its parameter budget (\S\ref{sec:impl-arch}); the full design is given
in \S\ref{sec:traj_enc}.

%------------------------------------------------------------------------------
\subsection{Trajectory Encoder Architecture (Brief)}
\label{sec:impl-arch}
The encoder has \textbf{35.80\,M} total parameters of which
\textbf{13.74\,M} are trainable; the frozen DINOv2-S/14 visual
backbone~\citep{oquab2023dinov2} (${\sim}22.1$\,M params, patch size $14$,
embedding dim $384$) accounts for the difference. Modality-specific widths
are $d_\text{traj}{=}128$, $d_\text{enc}{=}256$, $d_\text{vis}{=}256$, and
$d_\text{query}{=}128$. Component breakdown:

\begin{itemize}\itemsep1pt
\item \textbf{Query encoder} (1.46\,M): a $2$-layer, $4$-head transformer
      over tokenized question text with sinusoidal positional encoding,
      pooled to a single $128$-d query embedding.
\item \textbf{Visual patch encoder} (22.16\,M total; \textbf{0.10\,M}
      trainable): frozen DINOv2-S/14 over $16$ keyframes resized to
      $196{\times}196$ produces $14{\times}14{=}196$ patch tokens at $384$\,d,
      followed by a trainable linear projection to $d_\text{vis}{=}256$ and
      a temporal linear interpolation that aligns to the $T$ trajectory
      frames.
\item \textbf{Spatiotemporal trajectory encoder} (5.56\,M):
      a per-frame trajectory tokenizer producing four tokens
      $\{z^g, z^L, z^R, z^\phi\}$ (gaze, left hand, right hand, interaction),
      an intra-frame transformer (\emph{L1}: $2$ layers, $4$ heads,
      $d_\text{traj}{=}128$, FFN $4d$), an inter-frame temporal transformer
      (\emph{L2}: $6$ layers, $8$ heads, $d_\text{enc}{=}256$, FFN $4d$,
      sinusoidal positional encoding), a FiLM~\citep{perez2018film}
      block conditioned on the query embedding, and a
      visual--trajectory cross-attention fusion layer.
\item \textbf{Trajectory decoder} (3.20\,M): $3$-layer, $8$-head
      cross-attention decoder over future-frame queries, projecting to a
      $6$-d output (gaze, left-hand, right-hand 2D positions; sigmoid).
\item \textbf{Score decoder} (3.31\,M): a parallel $3$-layer, $8$-head
      cross-attention decoder projecting to $196$-d patch scores per
      future frame; used as auxiliary supervision in Stage~1 only.
\item \textbf{TrajScoreHead} (0.12\,M): an attention pool over the four
      per-frame trajectory tokens followed by a two-layer MLP
      (LayerNorm$\to$Linear$\to$GELU$\to$Linear) with a sigmoid head
      emitting the $196$-d patch importance map used at inference time
      and consumed by the Stage~2 token merger.
\end{itemize}

%------------------------------------------------------------------------------
\subsection{Stage 1: Trajectory-Driven Patch-Importance Pretraining}
\label{sec:impl-stage1}

\paragraph{Data.}
Stage~1 uses StreamGaze\_v2 clips drawn from
EgoExoLearn~\citep{huang2024egoexolearn} and
HoloAssist~\citep{wang2023holoassistegocentrichumaninteraction}, and
consumes only behavioral signals (gaze and hand trajectories) and image
frames; no question--answer labels are used. Each clip is uniformly sampled
into $T{=}128$ frames and split into a \emph{past} prefix of length $T_p$
and a \emph{future} suffix of length $T-T_p$, where the split ratio
$T_p/T$ is drawn uniformly from $[0.4, 0.6]$ at every training step. This
forces the model to support arbitrary $T_p$ at inference time. Gaze and
hand positions are loaded in normalized image coordinates with binary
visibility flags; missing modalities are handled by dedicated learnable
\emph{missing} embeddings inside the trajectory tokenizer.

\paragraph{Per-frame supervision targets.}
For every frame $t$ we render a ground-truth $14{\times}14$ patch-importance
map $\mathbf{I}_t \in [0,1]^{196}$ by summing isotropic Gaussian responses
on the patch-grid centers: a gaze Gaussian with bandwidth $\sigma_g{=}16$\,px
(in $224^2$ frame coordinates) plus the per-hand maximum of two hand
Gaussians with bandwidth $\sigma_h{=}24$\,px, with each modality contributing
zero whenever its visibility flag is off. The result is normalized per
frame and split into $\mathbf{I}^{\text{past}}_{1:T_p}$ and
$\mathbf{I}^{\text{future}}_{T_p+1:T}$.

\paragraph{Objective.}
The trajectory encoder is trained end-to-end with a sum of four equally
weighted MSE terms:
\begin{equation}
    \mathcal{L}_{\text{stage1}}
    \;=\; \mathcal{L}_\text{traj}
       \;+\; \mathcal{L}_\text{score-fut}
       \;+\; \mathcal{L}_\text{score-past}
       \;+\; \mathcal{L}_\text{score-traj} ,
    \label{eq:stage1}
\end{equation}
\begin{itemize}\itemsep1pt
\item $\mathcal{L}_\text{traj}$: future-trajectory regression in 2D image
      coordinates from the trajectory decoder, valid-length-masked over
      $T-T_p$ frames; gaze, left-hand, and right-hand positions are
      stacked into a $6$-d target.
\item $\mathcal{L}_\text{score-fut}$: future per-frame patch-score
      regression from the score decoder against
      $\mathbf{I}^{\text{future}}$.
\item $\mathcal{L}_\text{score-past}$: regression of the encoder's raw
      visual cross-attention readouts against $\mathbf{I}^{\text{past}}$,
      acting as a spatial-grounding auxiliary on the past frames.
\item $\mathcal{L}_\text{score-traj}$: distillation of the
      \emph{trajectory-driven} attention chain
      $\mathbf{A}^{\text{traj}}_{t,p}
      = \sum_{t'} \mathrm{decAttn}_{t', t}\!\cdot\!\mathrm{encAttn}_{t, p}$
      into the inference-time \texttt{TrajScoreHead}, with the chain
      target detached and per-frame max-normalized to $[0,1]$. This is
      the term that produces the patch scores consumed by Stage~2.
\end{itemize}
All four terms are computed with the variable-length valid mask induced by
$(T_p, T-T_p)$.

\paragraph{Optimization.}
Stage~1 is trained for \textbf{100 epochs} with
AdamW~\citep{loshchilov2019adamw}
($\beta_1{=}0.9$, $\beta_2{=}0.999$, weight decay $10^{-4}$), peak learning
rate $3{\times}10^{-4}$, and a cosine-annealing schedule decaying to
$3{\times}10^{-6}$ ($1\%$ of peak). The per-GPU batch size is $2$
(memory-bounded by the $T{=}128$ frame tensor); we train on
\textbf{four NVIDIA H200 GPUs under DistributedDataParallel} for an
effective global batch size of $8$. Frames are loaded at $224^2$ and
internally resized to $196^2$ on the DINOv2 path so that the $14{\times}14$
patch grid aligns exactly with $\mathbf{I}_t$. The checkpoint with the
lowest mean training loss is retained as the Stage~2 initialization.

%------------------------------------------------------------------------------
\subsection{Stage 2: Joint Fine-Tuning with VLM and Token Merging}
\label{sec:impl-stage2}

\paragraph{Vision--language backbone and adapters.}
We use \textbf{Qwen2.5-VL-7B-Instruct}~\citep{bai2025qwen25vl} as the
frozen vision--language backbone, loaded in \texttt{bfloat16}. Both the
visual encoder and the language model are kept frozen; only LoRA
adapters~\citep{hu2021loralowrankadaptationlarge} are trained on the LLM
side. LoRA is injected \emph{only} into the four self-attention projections
$\{q_\text{proj}, k_\text{proj}, v_\text{proj}, o_\text{proj}\}$ of every
LLM transformer layer (visual-encoder attention is not adapted), with rank
$r{=}16$, scaling $\alpha{=}32$, dropout $0.05$, and no bias term. This
yields $10{,}092{,}544$ trainable LoRA parameters out of
$8{,}302{,}259{,}200$ ($0.122\%$). The Stage~1 trajectory encoder is loaded
into the same training graph and fine-tuned end-to-end alongside the LoRA
adapters.

\paragraph{VLM input prompt.}
For every multiple-choice question we construct a single chat-formatted
\texttt{user} message with one video block followed by one text block. The
text uses the fixed template
\begin{quote}\small
\texttt{You are watching a short first-person (egocentric) video clip.}\\
\texttt{Question: \{q\}}\\\\
\texttt{\{options\}}\\\\
\texttt{Answer with only the letter (A, B, C, or D) of the correct option.}
\end{quote}
where \texttt{\{options\}} lists the four candidate answers as
\texttt{A.\,$\cdot$\,/\,B.\,$\cdot$\,/\,C.\,$\cdot$\,/\,D.\,$\cdot$}, one
per line. Frames are passed via Qwen's chat template with explicit
\texttt{resized\_height}~$=$~\texttt{resized\_width}~$=224$. Predictions are
obtained from a \emph{single} forward pass: we read the next-token logits at
the last prompt position and take the $\argmax$ over the four token IDs
corresponding to the letters $\{\text{A},\text{B},\text{C},\text{D}\}$. No
generation is performed and no length normalization is applied, so the task
becomes a four-way classification on top of the LM head.

\paragraph{Frame sampling and visual-token counts.}
Each clip is uniformly sampled into $T{=}128$ RGB frames for the VLM input
and a parallel $T_p{=}128$ frames for the trajectory encoder; frames are
resized to $224^2$. With Qwen2.5-VL's spatial-merge size $2$, temporal
patch size $2$, and patch size $14$, every $224^2$ frame contributes
$(224/14/2)^2{=}64$ spatial tokens, and the $128$ frames are folded into
$T_\text{merged}{=}64$ temporal groups, giving
\begin{equation}
    N \;=\; T_\text{merged}\cdot n_\text{spatial}
         \;=\; 64 \times 64 \;=\; 4{,}096
    \label{eq:n_video}
\end{equation}
visual tokens per sample. The text prompt contributes roughly $80$--$110$
tokens, so the full-resolution sequence length fed to the LLM is
$L_\text{full}\approx 4{,}200$.

\paragraph{Token reduction and compute savings.}
Let $\rho \in (0,1)$ denote the merge ratio (the fraction of visual tokens
removed) and $r{=}\lfloor \rho N \rfloor$ the number of removed tokens.
Bipartite gaze-weighted merging shortens the LLM input from $L_\text{full}$
to $L_\rho \approx (1-\rho)\,N + L_\text{text}$. Because the visual encoder
is frozen and the merge operates between the ViT and the LLM, ViT FLOPs are
unchanged and the savings accrue to the language model. Per LLM layer,
self-attention scales as $\mathcal{O}(L^2 d)$ and the MLP as
$\mathcal{O}(L d^2)$, with $d{=}3584$ for Qwen2.5-VL-7B.
Table~\ref{tab:compute} reports the per-layer FLOP reduction.

\begin{table}[t]
  \centering
  \small
  \caption{LLM compute as a function of the merge ratio $\rho$. Per-layer
  FLOPs combine self-attention ($\propto L^2 d$) and MLP
  ($\propto L d^2$); ViT cost is unchanged. The reduction factor is
  relative to the no-merge baseline. All FLOP and wall-clock measurements
  are taken on a single NVIDIA H200 GPU in training-mode
  forward$+$backward.}
  \label{tab:compute}
  \begin{tabular}{lcccc}
    \toprule
    Setting & Keep ratio & Kept tokens & LLM seq.\ length $L$
            & FLOPs reduction \\
    \midrule
    Full (no merge) & $100\%$  & $4{,}096$ & ${\sim}4{,}200$ & $1.0\times$  \\
    $\rho = 0.90$   & $10\%$   & $410$     & ${\sim}510$     & $15.6\times$ \\
    $\rho = 0.95$   & $5\%$    & $205$     & ${\sim}305$     & $27.5\times$ \\
    $\rho = 0.99$   & $1\%$    & $41$      & ${\sim}141$     & $62.2\times$ \\
    \bottomrule
  \end{tabular}
\end{table}

\noindent
At our default $\rho{=}0.95$ the LLM forward FLOPs drop by ${\sim}27.5\times$
and the KV cache shrinks correspondingly. Qwen2.5-VL-7B uses grouped-query
attention with $4$ KV heads and head dimension $128$, giving a KV feature
dimension of $4\times128{=}512$ per layer. A full-context cache therefore
requires $2\cdot 28 \cdot 4{,}200 \cdot 512 \cdot 2\,\text{B} \approx 241$\,MB
in \texttt{bfloat16}, whereas at $\rho{=}0.95$ it is ${\approx}18$\,MB.

\paragraph{Score alignment.}
The \texttt{TrajScoreHead} produces per-frame importance maps of shape
$(T_\text{traj}, 196)$, corresponding to a $14{\times}14$ patch grid at each of
the $T_\text{traj}{=}128$ trajectory frames. These must be aligned to the
$N{=}4{,}096$ visual tokens that Qwen2.5-VL presents to the LLM. The alignment
proceeds in two steps. \emph{Spatially}, each $14{\times}14$ map is
nearest-neighbor upsampled to $16{\times}16$ and then $2{\times}2$
average-pooled to $8{\times}8$, yielding $n_\text{spatial}{=}64$ scores per
frame—matching the $64$ spatial tokens per temporal group emitted by Qwen's
spatial-merge stride of $2$. (The intermediate upsample to $16{\times}16$ is
necessary because $14 \div 2 = 7 \neq 8$; direct half-resolution pooling would
yield a $7{\times}7$ grid that mismatches the VLM token layout.) \emph{Temporally}, the resulting
$(T_\text{traj}, 64)$ score tensor is linearly interpolated along the time axis
to the $T_\text{merged}{=}64$ temporal groups produced by Qwen's temporal-merge
stride of $2$. The aligned tensor is then flattened to
$\mathbf{s}\in\mathbb{R}^{N}$, giving one score per visual token.

\paragraph{Sequence construction with merged tokens.}
Given the per-token gaze-importance scores $\mathbf{s}\in\mathbb{R}^{N}$
emitted by the trajectory encoder, the top $(N-r)$ positions by score are
designated as \emph{receivers} and the bottom $r$ as \emph{sources}. Each
source is matched to the most cosine-similar receiver in token-embedding
space, and a score-weighted average is computed via \texttt{scatter\_add},
yielding $(N-r)$ merged embeddings. The source positions are then
\emph{removed} from the input sequence (rather than masked): the
corresponding entries of \texttt{input\_ids}, \texttt{attention\_mask}, and
the 3D-RoPE \texttt{position\_ids} are sliced out, and the remaining
receiver positions are populated with the merged embeddings via
\texttt{inputs\_embeds}. Both training and inference therefore operate at
the shortened length $L_\rho$, so the savings of Table~\ref{tab:compute}
apply identically to forward and backward passes.

\paragraph{Training procedure.}
We train Stage~2 for three epochs over the union of EgoExoLearn and
HoloAssist (5{,}799 items in 8 MCQ tasks; see \S\ref{sec:data}) and report
on EGTEA~\citep{li2018egtea} (526 items). Optimization uses
AdamW~\citep{loshchilov2019adamw} with weight decay $10^{-4}$, learning
rate $1{\times}10^{-4}$ for the LoRA adapters and $1{\times}10^{-5}$ for
the trajectory encoder, a constant schedule (no warmup), and global
gradient-norm clipping at $1.0$. The per-device batch size is fixed at one
and we accumulate gradients over four steps, yielding an effective batch
size of four. Stage~2 runs on a \textbf{single NVIDIA H200 GPU} in
\texttt{bfloat16}; deterministic seeding (\texttt{42}) is applied to
\texttt{torch}, \texttt{torch.cuda}, \texttt{random}, and \texttt{numpy}.
Hyperparameters are summarized in Table~\ref{tab:hparams}.

\begin{table}[t]
  \centering
  \small
  \caption{Stage-2 training hyperparameters.}
  \label{tab:hparams}
  \begin{tabular}{lll}
    \toprule
    Component & Hyperparameter & Value \\
    \midrule
    \multirow{4}{*}{Optimizer}
      & Optimizer            & AdamW \\
      & Weight decay         & $1\times10^{-4}$ \\
      & LoRA learning rate   & $1\times10^{-4}$ \\
      & TrajGaze encoder LR  & $1\times10^{-5}$ \\
    \midrule
    \multirow{3}{*}{Schedule}
      & Schedule             & constant, no warmup \\
      & Epochs               & 3 \\
      & Gradient clipping    & 1.0 (global $\ell_2$) \\
    \midrule
    \multirow{3}{*}{Batching}
      & Per-device batch     & 1 \\
      & Gradient accumulation & 4 \\
      & Effective batch      & 4 \\
    \midrule
    \multirow{3}{*}{LoRA}
      & Targets              & $\{q,k,v,o\}_\text{proj}$ in LLM only \\
      & Rank / $\alpha$ / dropout & $16$ / $32$ / $0.05$ \\
      & Trainable params     & $10{,}092{,}544$ ($0.122\%$) \\
    \midrule
    \multirow{4}{*}{Inputs}
      & VLM frames $T$              & 128 \\
      & Trajectory frames $T_p$     & 128 \\
      & Visual keyframes (DINOv2)   & 16 \\
      & Frame resolution            & $224\times224$ \\
    \midrule
    Hardware & GPU & 1$\times$ NVIDIA H200 (\texttt{bfloat16}) \\
    \bottomrule
  \end{tabular}
\end{table}

\paragraph{Training objective.}
Let $\boldsymbol{\ell}^{\text{stu}} \in \mathbb{R}^{|\mathcal{V}|}$ denote
the last-position LLM logits over the merged sequence and let
$y^\star \in \{\text{A},\text{B},\text{C},\text{D}\}$ be the gold answer
letter. We restrict the logits to the four option-letter token IDs,
$\boldsymbol{\ell}^{\text{stu}}_{\mathcal{C}} \in \mathbb{R}^{4}$, and apply
cross-entropy:
\begin{equation}
    \mathcal{L}_\text{CE}
    \;=\; -\log
    \frac{\exp\!\bigl(\boldsymbol{\ell}^{\text{stu}}_{\mathcal{C}}[y^\star]\bigr)}
         {\sum_{c\in\mathcal{C}} \exp\!\bigl(\boldsymbol{\ell}^{\text{stu}}_{\mathcal{C}}[c]\bigr)} .
    \label{eq:ce}
\end{equation}
For the no-distillation ablation reported in this paper,
$\mathcal{L} = \mathcal{L}_\text{CE}$. For the distilled variants
(\S\ref{sec:kd}) we additionally compute the same logits under a frozen
full-token teacher (Qwen2.5-VL-7B with the baseline-LoRA adapters), denoted
$\boldsymbol{\ell}^{\text{tea}}_{\mathcal{C}}$, and add a Kullback--Leibler
distillation term:
\begin{equation}
    \mathcal{L}
    \;=\; \alpha \cdot
    \mathrm{KL}\!\Bigl(
        \mathrm{softmax}(\boldsymbol{\ell}^{\text{tea}}_{\mathcal{C}})
        \,\big\|\,
        \mathrm{softmax}(\boldsymbol{\ell}^{\text{stu}}_{\mathcal{C}})
    \Bigr)
    + (1-\alpha)\cdot \mathcal{L}_\text{CE} .
    \label{eq:kd}
\end{equation}
Setting $\alpha{=}0$ recovers the no-KD trainer; gradients on the teacher
side are blocked at all times.

\paragraph{Evaluation protocol.}
Every $400$ optimizer steps we evaluate the entire EGTEA test split
(526 items) by computing top-1 accuracy under the same merged-token
inference pipeline used at training, and we keep the checkpoint that
maximizes this merged accuracy. Final per-task numbers are reported on the
selected checkpoint with $\rho$ fixed to its training value, with no
test-time data augmentation, generation, or temperature tuning.

%------------------------------------------------------------------------------
\subsection{Reproducibility}
\label{sec:impl-repro}
Code, configuration files, and exact launch scripts are released alongside
the paper. The Qwen2.5-VL-7B-Instruct snapshot hash used in all experiments
is \texttt{cc594898137f460bfe9f0759e9844b3ce807cfb5}.
Stage~1 uses four NVIDIA H200 GPUs under DDP and Stage~2 uses a
single NVIDIA H200 GPU; all training runs use a fixed random seed
(\texttt{42}) applied to \texttt{torch}, \texttt{torch.cuda},
\texttt{random}, and \texttt{numpy}. Stage~2 is launched without
distributed-data-parallel wrapping in order to eliminate cross-rank
floating-point variability and make single-GPU runs bit-stable for a given
seed.